\subsubsection{Measure Relationships}
\label{sec:measure_relationships}

Beyond postulate compliance, the following propositions establish how the three measures relate to each other. Rather than detecting independent failure modes, the measures exhibit a containment structure: sequencing conflicts are a strict subset of mutex relationships, while unreachability is independent of both.

\begin{prop}[Measure Relationships]
    \label{prop:measure_relationships}
    For any planning problem $\Pi$ and horizon $h$, the proposed measures exhibit the following relationships:
    \begin{enumerate}
        \item If goal $g$ is unreachable at horizon $h$, then $g$ does not participate in any
              mutex relationship or sequencing conflict at horizon $h$
        \item If $(g_1, g_2)$ is a sequencing conflict at horizon $h$ (with $g_2$ achievable),
              then $g_1$ and $g_2$ are mutex at horizon $h$
        \item Goals can be mutex at horizon $h$ without having any sequencing conflict at horizon $h$
    \end{enumerate}
\end{prop}
\begin{proof}
    \begin{enumerate}
        \item Both mutex (Definition~\ref{defn:mutex_goal_pair}) and sequencing conflict
              (Definition~\ref{defn:sequencing_conflict}) require participating goals to be achievable
              within $S^h_{\text{reach}}$.
              If $g$ is unreachable at horizon $h$, it fails this condition.

        \item Suppose $(g_1, g_2)$ is a sequencing conflict at horizon $h$ with $g_2$ achievable.
              By Definition~\ref{defn:sequencing_conflict}, $g_1$ is achievable and no execution trace
              of length at most $h$ achieves $g_1$ at time $i$ and $g_2$ at time $j \geq i$.
              In particular, no state in $S^h_{\text{reach}}$ contains both $g_1$ and $g_2$ (which would
              correspond to $i = j$). Combined with both goals being achievable,
              this satisfies Definition~\ref{defn:mutex_goal_pair}.

        \item Light Switch: Goals $\{\text{on}, \text{off}\}$ are mutex
              (both achievable, never co-occur) but have no sequencing conflict (from any
              \textit{on}-state, \textit{off} is reachable via \textit{turn\_off}, and vice versa).
    \end{enumerate}
\end{proof}

\begin{figure}[ht]
    \centering
    \begin{tikzpicture}[
            region/.style={rectangle, draw, thick, rounded corners=4pt},
            mlabel/.style={font=\small},
        ]
        % Sequencing label (right side, lower)
        \node[mlabel] (gs-label) at (1.3, -1)
        {\textbf{Sequencing} $\mathcal{I}_{\text{GS}}$};

        % Mutex label (right side, upper)
        \node[mlabel] (mx-label) at (0.7, 0)
        {\textbf{Mutex} $\mathcal{I}_{\text{MX}}$};

        % Unreachability label (left side, upper)
        \node[mlabel] (ur-label) at (-2.8, 0)
        {\textbf{Unreachability} $\mathcal{I}_{\text{UR}}$};

        % Sequencing box (tight around label)
        \node[region, fit=(gs-label), inner xsep=8pt, inner ysep=6pt] (gs) {};

        % Mutex box (around label + sequencing)
        \node[region, fit=(mx-label)(gs), inner xsep=10pt, inner ysep=6pt] (mx) {};

        % Unreachability box (match mutex height exactly)
        \node[region, fit=(ur-label)(ur-label |- mx.south)(ur-label |- mx.north),
            inner xsep=10pt, inner ysep=0pt] (ur) {};

        % Outer label (placed first, then included in fit so box expands to hold it)
        \node[font=\small\bfseries] (all-label) at (-0.7, 0.8)
        {All (Goal) Propositions};

        % Outer box (fits inner boxes AND the label)
        \node[region, fit=(ur)(mx)(all-label), inner xsep=12pt, inner ysep=6pt] (all) {};

        % --- Fills (background layer: outside to inside) ---
        \begin{scope}[on background layer]
            \fill[black!5, rounded corners=4pt]
            (ur.south west) rectangle (ur.north east);
            \fill[black!8, rounded corners=4pt]
            (mx.south west) rectangle (mx.north east);
            \fill[black!15, rounded corners=4pt]
            (gs.south west) rectangle (gs.north east);
        \end{scope}

    \end{tikzpicture}
    \caption[Diagnostic measure relationships.]{Diagnostic measure relationships. Sequencing conflicts are a subset of mutex relationships; unreachable (goal) propositions are excluded from both measures.}
    \label{fig:measure_relationships}
\end{figure}


The scope and structure variants of each measure are also formally related:

\begin{prop}[Scope-Structure Bounds]
    \label{prop:scope_structure_bounds}
    For any planning problem $\Pi$ and horizon $h$:
    \begin{enumerate}
        \item $\mathcal{I}^{\text{scope}}_{\text{UR}}(\Pi, h) \leq \mathcal{I}^{\text{struct}}_{\text{UR}}(\Pi, h) \leq |P|$
        \item $\mathcal{I}^{\text{struct}}_{\text{MX}}(\Pi, h) \leq \binom{\mathcal{I}^{\text{scope}}_{\text{MX}}(\Pi, h)}{2}$
        \item $\mathcal{I}^{\text{struct}}_{\text{GS}}(\Pi, h) \leq \mathcal{I}^{\text{scope}}_{\text{GS}}(\Pi, h) \times (\mathcal{I}^{\text{scope}}_{\text{GS}}(\Pi, h) - 1)$
    \end{enumerate}
\end{prop}
\begin{proof}
    \begin{enumerate}
        \item Unreachable goals are a subset of all unreachable propositions,
              which are a subset of $P$.

        \item The maximum number of unordered pairs among
              $\mathcal{I}^{\text{scope}}_{\text{MX}}$ goals is $\binom{\mathcal{I}^{\text{scope}}_{\text{MX}}}{2}$.

        \item The maximum number of ordered pairs among
              $\mathcal{I}^{\text{scope}}_{\text{GS}}$ goals is $n(n-1)$ where
              $n = \mathcal{I}^{\text{scope}}_{\text{GS}}$.
    \end{enumerate}
\end{proof}
